\section{Rynek mocy --- the Polish capacity market}

The day-ahead market pays for \emph{energy delivered}. The capacity
market pays for \emph{being available}. Poland runs one of Europe's
largest; it appears in nearly every Polish energy job description.

\subsection{Why it exists}

A system dominated by cheap renewables has hours where prices are near
zero. A gas peaker that runs 200 hours a year cannot recover its fixed
costs from energy sales alone --- the ``missing money'' problem. Without
extra revenue it retires, and then a calm, dark January evening has no
capacity left to call. The capacity market pays plants (and demand
response) a fixed fee for guaranteed availability in stress events.

\subsection{Mechanics, briefly}

\begin{itemize}
  \item Introduced by the 2017 capacity market act; first delivery year
        2021. Modeled on the British design, EU state-aid approved.
  \item \textbf{Descending-clock auctions} held years ahead (main
        auction $\sim$5 years before delivery, plus adjustment
        auctions). Price falls each round until supply matches the
        demand curve set by the regulator.
  \item Winners take on a \emph{capacity obligation}: be available
        during announced stress periods, or face penalties. New builds
        can lock contracts for up to 15--17 years; existing units get
        1 year.
  \item Paid for by consumers through the \emph{op\l{}ata mocowa}
        (capacity fee) on electricity bills.
  \item The EU 550g CO$_2$/kWh emission limit pushed unmodernized coal
        out of new contracts --- a major driver of the Polish coal fleet's
        retirement schedule, and therefore of long-run price formation.
\end{itemize}

\subsection{Why a forecasting desk cares}

\begin{itemize}
  \item Capacity revenues decide which plants stay in the merit order.
        The retirement schedule shifts the supply curve our price models
        implicitly learn --- a regime change models must re-learn.
  \item Stress-event risk shows up as scarcity premiums in the
        day-ahead and balancing prices in tight winter hours --- exactly
        the spike tail where our models are weakest (see the worst-days
        panel in \texttt{reports/figures/backtests/}).
  \item For a trader, announced capacity auction results are slow-moving
        fundamentals: they change the expected frequency of spike hours
        years ahead.
\end{itemize}

Interview line: ``Day-ahead pays for energy, rynek mocy pays for
insurance. The insurance side decides which plants exist, so it quietly
sets the shape of the price distribution my models learn.''
